% ============================================================
% 附录 A：Abel 求和与阶梯函数上的 Stieltjes 积分（正文第2章起）
% ============================================================
\chapter{Abel 求和与阶梯函数上的 Stieltjes 积分}
\label{app:abel_stieltjes}

本附录给出本书所用的\textbf{分部求和（Abel 求和）}与\textbf{阶梯函数上的 Stieltjes 积分}
的工作定义及必要推演。范围限于\textbf{右连续阶梯函数}（有限或可数跳跃），
足以覆盖第~\ref{ch:elementary_number_theoretic_functions}~章的
$\pi,J,\theta,\psi$、第~\ref{ch:integral_transforms}~章的 Mellin--Perron 准备，
以及第~\ref{ch:riemann_explicit_formula}~章的 $\psi\to J$ 分部。
\textbf{不}展开一般有界变差函数或测度论 Stieltjes 积分；
亦\textbf{不}讨论 Stieltjes \textbf{变换}（与此处积分不是同一对象）。

正文首次使用处（第~\ref{ch:elementary_number_theoretic_functions}~章）
只保留「跳跃和怎么算」的短句，系统表述以本附录为锚点。

%------------------------------------------------------------
\section{右连续阶梯函数}
\label{sec:app-step-functions}
%------------------------------------------------------------

\begin{definition}{}{}
称 $A:[a,\infty)\to\mathbb{R}$ 为\textbf{右连续阶梯函数}，若存在至多可数的
严格递增序列 $(x_i)$（跳跃点）与实数 $(a_i)$（跳跃量），使得
$A$ 在相邻跳跃点之间为常数，且
\[
  A(x_i)-A(x_i^-) \;=\; a_i,
  \qquad
  A(x_i) \;=\; \lim_{t\to x_i^+} A(t)
\]
（右连续：跳跃点处取跳跃后的值）。
\end{definition}

本书主要例子（均取 $A(x)=0$ 当 $x$ 小于第一个有贡献的点）：
\begin{center}
\renewcommand{\arraystretch}{1.35}
\begin{tabular}{@{}lll@{}}
\toprule
函数 & 跳跃位置 & 跳跃量 \\
\midrule
$\pi(x)$ & 素数 $p$ & $1$ \\
$J(x)$ & 素数幂 $p^k$ & $1/k$ \\
$\theta(x)$ & 素数 $p$ & $\log p$ \\
$\psi(x)$ & 素数幂 $p^k$ & $\log p$ \\
$\displaystyle\sum_{n\le x}a_n$ & 正整数 $n$ & $a_n$ \\
\bottomrule
\end{tabular}
\end{center}

%------------------------------------------------------------
\section{Abel 求和（分部求和）}
\label{sec:app-abel}
%------------------------------------------------------------

设 $a_n\in\mathbb{C}$，$A(x)=\sum_{n\le x}a_n$（右连续阶梯），
$f$ 在 $[1,x]$ 上连续可微。则有\textbf{分部求和公式}
\begin{equation}
  \sum_{1\le n\le x} a_n f(n)
  \;=\;
  A(x)f(x)
  \;-\;
  \int_1^{x} A(t)\,f'(t)\,\mathrm{d}t,
  \label{eq:app-abel}
\end{equation}
其中积分是通常的 Riemann 积分（$A$ 在有限个间断点处取值不影响积分）。
若求和从 $n=n_0$ 开始，则下限改为 $n_0^-$，$A$ 在起点左侧取左极限。

\begin{proof}[轮廓]
在每个区间 $[n,n+1)$ 上 $A$ 为常数 $A(n)$，故
\[
\int_n^{n+1} A(t)f'(t)\,\mathrm{d}t
  \;=\;
  A(n)\bigl(f(n+1)-f(n)\bigr).
\]
对 $n=1,\ldots,\lfloor x\rfloor-1$ 求和并与边界项 $A(x)f(x)-A(1^-)f(1)$ 合并，
利用 $A(n)-A(n-1)=a_n$ 即得 \eqref{eq:app-abel}。
（有限项时直接验证；无穷级数情形在绝对收敛或部分和有界等假设下取极限。）
\end{proof}

式 \eqref{eq:app-abel} 亦称 \textbf{Abel 求和}。它是离散部分和与光滑权重之间的基本桥梁。

%------------------------------------------------------------
\section{阶梯函数上的 Stieltjes 积分}
\label{sec:app-stieltjes}
%------------------------------------------------------------

\begin{definition}{阶梯上的 Stieltjes 积分（工作定义）}{app-stieltjes}
设 $A$ 为右连续阶梯函数，跳跃点 $x_i$、跳跃量 $a_i$；
$f$ 在每个 $x_i$ 处有定义（本书应用中 $f$ 连续即可）。
对 $a<b$，定义
\begin{equation}
  \int_{(a,b]} f(t)\,\mathrm{d}A(t)
  \;=\;
  \sum_{\substack{i \\ a < x_i \le b}} f(x_i)\,a_i.
  \label{eq:app-stieltjes-def}
\end{equation}
左端亦可写 $\int_{a^+}^{b}$ 或（在约定 $A$ 于 $a$ 左侧无质量时）$\int_a^b$。
本书常用 $\int_{2^-}^{x}$ 表示自 $2$ 的左邻域积到 $x$（含 $t=2$ 处的跳跃）。
\end{definition}

因而：对阶梯 $A$，Stieltjes 积分\textbf{就是}「$f$ 在跳跃点的值 × 跳跃量」的和。
特别地，若 $A(x)=\sum_{n\le x}a_n$，则
\begin{equation}
  \int_{(0,x]} f(t)\,\mathrm{d}A(t)
  \;=\;
  \sum_{n\le x} f(n)\,a_n.
  \label{eq:app-stieltjes-sum}
\end{equation}

\begin{remark}{与 Stieltjes 变换的区别}{}
\textbf{Stieltjes 变换}指形如 $\int (z-t)^{-1}\,\mathrm{d}\mu(t)$ 的解析对象；
本附录的 $\int f\,\mathrm{d}A$ 是\textbf{积分}本身，二者不可混称。
\end{remark}

%------------------------------------------------------------
\section{分部积分公式}
\label{sec:app-stieltjes-parts}
%------------------------------------------------------------

\begin{theorem}{阶梯 Stieltjes 分部积分}{app-stieltjes-parts}
设 $A$ 右连续阶梯，$f\in C^1[a,b]$。则
\begin{equation}
  \int_a^b f(t)\,\mathrm{d}A(t)
  \;=\;
  f(b)A(b)
  \;-\;
  f(a)A(a^-)
  \;-\;
  \int_a^b A(t)\,f'(t)\,\mathrm{d}t,
  \label{eq:app-stieltjes-parts}
\end{equation}
其中 $A(a^-)=\lim_{t\to a^-}A(t)$；若 $a$ 小于一切跳跃点，则 $A(a^-)=A(a)$ 或取 $0$
（依 $A$ 在 $a$ 左侧的约定而定）。
\end{theorem}

\begin{proof}
左端按定义 \eqref{eq:app-stieltjes-def} 为 $\sum_{a<x_i\le b}f(x_i)a_i$。
右端边界项展开跳跃后，与 Abel 公式 \eqref{eq:app-abel} 在「$a_n$ 仅在跳跃点非零」时一致；
亦可在每个连续区间上对常数 $A$ 作普通分部再把跳跃贡献单独相加。
\end{proof}

因而：\textbf{Abel 求和是离散版；阶梯 Stieltjes 分部是同一思想的积分写法。}

%------------------------------------------------------------
\section{本书中的标准应用}
\label{sec:app-abel-stieltjes-apps}
%------------------------------------------------------------

下列恒等式在第~\ref{ch:elementary_number_theoretic_functions}~章给出数论含义；
此处仅标明它们如何落入 \eqref{eq:app-stieltjes-def}、\eqref{eq:app-stieltjes-parts}。

\subsection{\texorpdfstring{$\theta$}{theta} 与 \texorpdfstring{$\pi$}{pi}}
$\mathrm{d}\pi$ 在素数 $p$ 处质量为 $1$，故
\[
\theta(x)
  \;=\;
  \sum_{p\le x}\log p
  \;=\;
  \int_{2^-}^{x}\log t\,\mathrm{d}\pi(t).
\]
对 $f(t)=\log t$ 用 \eqref{eq:app-stieltjes-parts}（或直接 Abel）得
\begin{equation}
  \theta(x)
  \;=\;
  \pi(x)\log x
  \;-\;
  \int_2^{x}\frac{\pi(t)}{t}\,\mathrm{d}t
  \label{eq:app-theta-pi}
\end{equation}
（与正文 \eqref{eq:theta_pi_abel} 相同）。

\subsection{\texorpdfstring{$\psi$}{psi} 与 \texorpdfstring{$J$}{J}}
$\mathrm{d}J$ 在 $p^k$ 处质量 $1/k$，$\mathrm{d}\psi$ 在 $p^k$ 处质量 $\log p$，故
\begin{align}
\psi(x)
  &=
  \int_{2^-}^{x}\log t\,\mathrm{d}J(t),
  \label{eq:app-psi-from-J}\\
J(x)
  &=
  \int_{2^-}^{x}\frac{1}{\log t}\,\mathrm{d}\psi(t)
  =
  \sum_{2\le n\le x}\frac{\Lambda(n)}{\log n}.
  \label{eq:app-J-from-psi}
\end{align}
对 \eqref{eq:app-psi-from-J} 用 \eqref{eq:app-stieltjes-parts} 得
\begin{equation}
  \psi(x)
  \;=\;
  \log x\cdot J(x)
  \;-\;
  \int_1^{x}\frac{J(t)}{t}\,\mathrm{d}t.
  \label{eq:app-psi-J-parts}
\end{equation}
对 \eqref{eq:app-J-from-psi} 用 $f=1/\log t$ 得第~\ref{ch:riemann_explicit_formula}~章所用形式
\begin{equation}
  J(x)
  \;=\;
  \frac{\psi(x)}{\log x}
  \;+\;
  \int_2^{x}\frac{\psi(t)}{t(\log t)^2}\,\mathrm{d}t
  \qquad (x>1),
  \label{eq:app-J-from-psi-parts}
\end{equation}
（因 $\psi(2^-)=0$）。详细展开见该章第~\ref{sec:integration-by-parts}~节；
一般公式以本附录 \eqref{eq:app-stieltjes-parts} 为准。

\subsection{Mellin 与部分和}
若 $A(x)=\sum_{n\le x}a_n$，$F(s)=\sum a_n n^{-s}$，则在收敛半平面内
\[
\int_1^{\infty} A(x)\,x^{-s-1}\,\mathrm{d}x
  \;=\;
  \frac{F(s)}{s}
\]
（对 $A$ 与幂函数作 Abel / 分部；见第~\ref{ch:integral_transforms}~章）。

%------------------------------------------------------------
\section{小结}
\label{sec:app-abel-stieltjes-summary}
%------------------------------------------------------------

\begin{center}
\renewcommand{\arraystretch}{1.45}
\begin{tabular}{@{}p{3.2cm}p{10cm}@{}}
\toprule
\textbf{对象} & \textbf{本书用法} \\
\midrule
Abel 求和 \eqref{eq:app-abel} &
  离散 $a_n$ 与光滑 $f$；$\theta$–$\pi$、Dirichlet 级数与部分和 \\
Stieltjes（阶梯）\eqref{eq:app-stieltjes-def} &
  $\int f\,\mathrm{d}A=\sum f(x_i)a_i$；$\psi$–$J$、Perron 准备 \\
分部 \eqref{eq:app-stieltjes-parts} &
  显式公式中 $\psi\to J$、振幅 $1/\log x$ 的来源 \\
\bottomrule
\end{tabular}
\end{center}

\vspace{0.5em}
\noindent
\textbf{正文指针：}
第~\ref{ch:elementary_number_theoretic_functions}~章第~\ref{subsec:theta_pi_abel}~、
\ref{subsec:J_psi_stieltjes}~节（应用）；
第~\ref{ch:integral_transforms}~章（Mellin / Perron 语境）；
第~\ref{ch:riemann_explicit_formula}~章第~\ref{sec:integration-by-parts}~节（$\psi\to J$ 展开）；
第~\ref{ch:fourier_vs_riemann}~章（$\psi\to J\to\pi$ 与 Fourier 比较中的 Abel 步骤）。
